\documentclass[11pt,a4paper]{article}
\usepackage[utf8]{inputenc}
\usepackage[T1]{fontenc}
\usepackage[margin=1in]{geometry}
\begin{document}








\maketitle


\part*{Part I: Literature Review}\addcontentsline{toc}{part}{Part I: Literature Review}

\documentclass[11pt,a4paper]{article}












% Define custom colors
\definecolor{promptblue}{RGB}{41, 98, 150}
\definecolor{promptbg}{RGB}{240, 247, 255}
\definecolor{promptborder}{RGB}{100, 160, 220}
\definecolor{accentgreen}{RGB}{0, 128, 100}
\definecolor{warningorange}{RGB}{200, 120, 0}

\newtcolorbox{resultbox}[1][]{
    colback=promptbg,
    colframe=promptborder,
    coltitle=white,
    fonttitle=\bfseries\small,
    title={#1},
    colbacktitle=promptblue,
    boxrule=1pt,
    arc=2mm,
    left=8pt, right=8pt, top=6pt, bottom=6pt,
    enhanced, breakable
}

\newtcolorbox{gapbox}[1][]{
    colback=orange!5,
    colframe=warningorange,
    coltitle=white,
    fonttitle=\bfseries\small,
    title={#1},
    colbacktitle=warningorange,
    boxrule=1pt,
    arc=2mm,
    left=8pt, right=8pt, top=6pt, bottom=6pt,
    enhanced, breakable
}

\hypersetup{
    colorlinks=true,
    linkcolor=promptblue,
    urlcolor=promptblue,
    citecolor=promptblue
}

\title{\textbf{Literature Review: Femomenología del color
}}
\author{Generated by Research Accelerant Agent}
\date{\today}

\begin{document}
\maketitle



\begin{resultbox}[Search Parameters]
\begin{itemize}[leftmargin=*]
    \item \textbf{Topic:} Femomenología del color

    \item \textbf{Year Range:} 2010--2026
    \item \textbf{Studies Retrieved:} 5
    \item \textbf{Citation Format:} APA
\end{itemize}
\end{resultbox}

\vspace{1em}

\textbf{\large Reviewed Studies}\par\vspace{0.5em}


\textbf{1. Teorías de la luz y el color en la época de las Luces. De Newton a Goethe}

\textbf{Full Citation:} J Pimentel (2015). Teorías de la luz y el color en la época de las Luces. De Newton a Goethe. 	extit{Arbor}. https://doi.org/N/A

\begin{resultbox}[Study Details]
\begin{itemize}[leftmargin=*]
    \item \textbf{Authors:} J Pimentel
    \item \textbf{Year:} 2015
    \item \textbf{Journal:} Arbor
    
    
    
    \item \textbf{DOI:} \href{https://doi.org/}{N/A}
    \item \textbf{Citation Count:} 29 (Google Scholar)
    \item \textbf{Access:} \href{https://arbor.revistas.csic.es/index.php/arbor/article/view/2067}{Semantic Scholar}
\end{itemize}
\end{resultbox}

\textbf{Abstract:} … por revisar la teoría de la luz de Isaac Newton, centrándonos … entendida la luz, así como el papel simbólico de la luz, metáfora … entender mejor no sólo qué son la luz, el ojo o los fenóme…\par\vspace{0.5em}



\textbf{2. As cores ea Visão ea Visão das Cores}

\textbf{Full Citation:} MCS Ribeiro (2011). As cores ea Visão ea Visão das Cores. 	extit{2011}. https://doi.org/N/A

\begin{resultbox}[Study Details]
\begin{itemize}[leftmargin=*]
    \item \textbf{Authors:} MCS Ribeiro
    \item \textbf{Year:} 2011
    \item \textbf{Journal:} 2011
    
    
    
    \item \textbf{DOI:} \href{https://doi.org/}{N/A}
    \item \textbf{Citation Count:} 10 (Google Scholar)
    \item \textbf{Access:} \href{https://search.proquest.com/openview/d00872f699d42bae8e5936ff65e42d0a/1?pq-origsite=gscholar&cbl=2026366&diss=y}{Semantic Scholar}
\end{itemize}
\end{resultbox}

\textbf{Abstract:} … A luz é considerada como a energia que sustenta a vida. Toda … Sabemos que através da luz recebemos grandes quantidades de … aos diferentes comprimentos de onda da luz recebida. …\par\vspace{0.5em}



\textbf{3. Evaluación de las estrategias de adaptación a disfunciones de la visión del color}

\textbf{Full Citation:} I Coca Torrents (2012). Evaluación de las estrategias de adaptación a disfunciones de la visión del color. 	extit{2012}. https://doi.org/N/A

\begin{resultbox}[Study Details]
\begin{itemize}[leftmargin=*]
    \item \textbf{Authors:} I Coca Torrents
    \item \textbf{Year:} 2012
    \item \textbf{Journal:} 2012
    
    
    
    \item \textbf{DOI:} \href{https://doi.org/}{N/A}
    \item \textbf{Citation Count:} 3 (Google Scholar)
    \item \textbf{Access:} \href{https://upcommons.upc.edu/entities/publication/286e34af-fb41-42b1-8d35-5fc558f4e080}{Semantic Scholar}
\end{itemize}
\end{resultbox}

\textbf{Abstract:} … nos ayuda, por ejemplo, a saber si una fruta está lo suficientemente madura para poder comerla, aun sabiendo que el color de dicha fruta puede cambiar dependiendo de la luz que le …\par\vspace{0.5em}



\textbf{4. La visión del color: mecanismos fisiológicos, teorías y deficiencias comunes}

\textbf{Full Citation:} MS Larraz, AC Martínez, BM Lacámara… (2025). La visión del color: mecanismos fisiológicos, teorías y deficiencias comunes. 	extit{Revista Sanitaria de …}. https://doi.org/N/A

\begin{resultbox}[Study Details]
\begin{itemize}[leftmargin=*]
    \item \textbf{Authors:} MS Larraz, AC Martínez, BM Lacámara…
    \item \textbf{Year:} 2025
    \item \textbf{Journal:} Revista Sanitaria de …
    
    
    
    \item \textbf{DOI:} \href{https://doi.org/}{N/A}
    \item \textbf{Citation Count:} 0 (Google Scholar)
    \item \textbf{Access:} \href{https://dialnet.unirioja.es/servlet/articulo?codigo=10349656}{Semantic Scholar}
\end{itemize}
\end{resultbox}

\textbf{Abstract:} … en color) y bastones (para la visión nocturna). Los conos están distribuidos en la fóvea y responden a las longitudes de onda de luz … In conclusion, color vision is fundamental both in …\par\vspace{0.5em}



\textbf{5. LA PERCEPCIÓN DE LA LUZ POR EL SISTEMA VISUAL HUMANO. EL COLOR}

\textbf{Full Citation:} JR Mora (2026). LA PERCEPCIÓN DE LA LUZ POR EL SISTEMA VISUAL HUMANO. EL COLOR. 	extit{Luz y Vida.}. https://doi.org/N/A

\begin{resultbox}[Study Details]
\begin{itemize}[leftmargin=*]
    \item \textbf{Authors:} JR Mora
    \item \textbf{Year:} 2026
    \item \textbf{Journal:} Luz y Vida.
    
    
    
    \item \textbf{DOI:} \href{https://doi.org/}{N/A}
    \item \textbf{Citation Count:} 0 (Google Scholar)
    \item \textbf{Access:} \href{https://www.rasc.es/wp-content/uploads/2025/01/Luz-y-Vida_Conmemorando-el-dia-internacional-de-la-luz.pdf#page=89}{Semantic Scholar}
\end{itemize}
\end{resultbox}

\textbf{Abstract:} … absoluta de luz y la saturación relacionada con la pureza del color que presente la luz en … sistema de representación del color que ayudaran a medir el color con números. Esto era …\par\vspace{0.5em}



\vspace{1em}
\textbf{\large Summary Statistics}\par\vspace{0.5em}

\begin{center}
\begin{tabular}{lc}
\toprule
\textbf{Metric} & \textbf{Value} \\
\midrule
Total Studies Reviewed & 5 \\
Average Citations per Study & 8 \\
Publication Year Range & 2011--2026 \\
Studies with Open Access PDF & 4 \\
\bottomrule
\end{tabular}
\end{center}


\vfill
\begin{center}
\textit{Document generated by Research Accelerant Agent on \today.}
\end{center}

\end{document}

\newpage

\part*{Part II: Cross-Study Synthesis}\addcontentsline{toc}{part}{Part II: Cross-Study Synthesis}

\documentclass[11pt,a4paper]{article}












% Define custom colors
\definecolor{promptblue}{RGB}{41, 98, 150}
\definecolor{promptbg}{RGB}{240, 247, 255}
\definecolor{promptborder}{RGB}{100, 160, 220}
\definecolor{accentgreen}{RGB}{0, 128, 100}
\definecolor{warningorange}{RGB}{200, 120, 0}

\newtcolorbox{resultbox}[1][]{
    colback=promptbg,
    colframe=promptborder,
    coltitle=white,
    fonttitle=\bfseries\small,
    title={#1},
    colbacktitle=promptblue,
    boxrule=1pt,
    arc=2mm,
    left=8pt, right=8pt, top=6pt, bottom=6pt,
    enhanced, breakable
}

\newtcolorbox{gapbox}[1][]{
    colback=orange!5,
    colframe=warningorange,
    coltitle=white,
    fonttitle=\bfseries\small,
    title={#1},
    colbacktitle=warningorange,
    boxrule=1pt,
    arc=2mm,
    left=8pt, right=8pt, top=6pt, bottom=6pt,
    enhanced, breakable
}

\hypersetup{
    colorlinks=true,
    linkcolor=promptblue,
    urlcolor=promptblue,
    citecolor=promptblue
}

\title{\textbf{Cross-Study Synthesis: Femomenología del color
}}
\author{Generated by Research Accelerant Agent}
\date{\today}

\begin{document}
\maketitle



\textbf{\large Studies Included in Synthesis}\par\vspace{0.5em}

\begin{longtable}{p{5cm}cc}
\toprule
\textbf{Study} & \textbf{Year} & \textbf{Citations} \\
\midrule
\endhead
Teorías de la luz y el color en la época de las Lu... & 2015 & 29 \\
As cores ea Visão ea Visão das Cores & 2011 & 10 \\
Evaluación de las estrategias de adaptación a disf... & 2012 & 3 \\
La visión del color: mecanismos fisiológicos, teor... & 2025 & 0 \\
LA PERCEPCIÓN DE LA LUZ POR EL SISTEMA VISUAL HUMA... & 2026 & 0 \\
\bottomrule
\end{longtable}

\vspace{1em}


\textbf{\large 1. Methodological Patterns}\par\vspace{0.5em}

Across the 5 reviewed studies on Femomenología del color
, the following methodological patterns emerged:

- **Not specified**: 5 studies

**Design Distribution:**
- Total experimental studies: 0
- Total meta-analyses/systematic reviews: 0
- Total longitudinal studies: 0
- Total qualitative studies: 0
- Total computational studies: 0

\vspace{1em}


\textbf{\large 2. Overarching Findings and Trends}\par\vspace{0.5em}

**Consensus Themes:**
- Multiple studies focus on: color, teor, cores, visi.

**Emergent Trends:**
- 2 recent studies (2025-2026) indicate evolving methodological approaches.

**Contradictions/Debates:**
- No major contradictions noted.

\vspace{1em}


\textbf{\large 3. Recurring Gaps and Limitations}\par\vspace{0.5em}

**Limitations noted by authors:**
- Authors of reviewed studies did not explicitly identify methodological limitations.

**Systematic gaps in the literature:**
1. Geographic scope: Most studies on Femomenología del color
 may be concentrated in specific regions, limiting generalizability to diverse populations.
2. Methodological narrowness: The literature relies heavily on Not specified approaches, with limited triangulation.
3. Sample and temporal constraints: Many studies use cross-sectional designs or limited sample sizes, limiting causal inference about Femomenología del color
.
4. Translation to practice: Limited evidence on how findings regarding Femomenología del color
 translate into actionable interventions or policy recommendations.

\vspace{0.5em}


\begin{gapbox}[Gap 1]
Geographic scope: Most studies on Femomenología del color
 may be concentrated in specific regions, limiting generalizability to diverse populations.
\end{gapbox}

\vspace{0.3em}


\begin{gapbox}[Gap 2]
Methodological narrowness: The literature relies heavily on Not specified approaches, with limited triangulation.
\end{gapbox}

\vspace{0.3em}


\begin{gapbox}[Gap 3]
Sample and temporal constraints: Many studies use cross-sectional designs or limited sample sizes, limiting causal inference about Femomenología del color
.
\end{gapbox}

\vspace{0.3em}


\begin{gapbox}[Gap 4]
Translation to practice: Limited evidence on how findings regarding Femomenología del color
 translate into actionable interventions or policy recommendations.
\end{gapbox}

\vspace{0.3em}


\vspace{1em}
\textbf{\large 4. Impact Assessment}\par\vspace{0.5em}

**High-Impact Studies (>$50$ citations):** 0 out of 5 reviewed studies (0%).

**Recent Contributions (2024+):** 2 studies.

**Most Cited Study:** Teorías de la luz y el color en la época de las Luces. De Newton a Goethe (29 citations).

**Open Access Availability:** 4 of 5 studies (80%) provide open access to their full text.

\vspace{1em}


\textbf{\large 5. Future Research Directions}\par\vspace{0.5em}

Based on the identified gaps, the following research directions are recommended:

1. **Geographic scope**: Further investigation needed to address geographic scope: most studies on femomenología del color
 may be concentrated in specific regions, limiting generalizability to diverse populations..
2. **Methodological narrowness**: Further investigation needed to address methodological narrowness: the literature relies heavily on not specified approaches, with limited triangulation..
3. **Sample and temporal constraints**: Further investigation needed to address sample and temporal constraints: many studies use cross-sectional designs or limited sample sizes, limiting causal inference about femomenología del color
..
4. **Translation to practice**: Further investigation needed to address translation to practice: limited evidence on how findings regarding femomenología del color
 translate into actionable interventions or policy recommendations..

\vspace{1em}


\vfill
\begin{center}
\textit{Synthesis generated by Research Accelerant Agent on \today.}
\end{center}

\end{document}

\newpage

\part*{Part III: Problem Statement}\addcontentsline{toc}{part}{Part III: Problem Statement}

\documentclass[11pt,a4paper]{article}












% Define custom colors
\definecolor{promptblue}{RGB}{41, 98, 150}
\definecolor{promptbg}{RGB}{240, 247, 255}
\definecolor{promptborder}{RGB}{100, 160, 220}
\definecolor{accentgreen}{RGB}{0, 128, 100}
\definecolor{warningorange}{RGB}{200, 120, 0}

\newtcolorbox{resultbox}[1][]{
    colback=promptbg,
    colframe=promptborder,
    coltitle=white,
    fonttitle=\bfseries\small,
    title={#1},
    colbacktitle=promptblue,
    boxrule=1pt,
    arc=2mm,
    left=8pt, right=8pt, top=6pt, bottom=6pt,
    enhanced, breakable
}

\newtcolorbox{gapbox}[1][]{
    colback=orange!5,
    colframe=warningorange,
    coltitle=white,
    fonttitle=\bfseries\small,
    title={#1},
    colbacktitle=warningorange,
    boxrule=1pt,
    arc=2mm,
    left=8pt, right=8pt, top=6pt, bottom=6pt,
    enhanced, breakable
}

\hypersetup{
    colorlinks=true,
    linkcolor=promptblue,
    urlcolor=promptblue,
    citecolor=promptblue
}

\title{\textbf{Problem Statement: Femomenología del color
}}
\author{Generated by Research Accelerant Agent}
\date{\today}

\begin{document}
\maketitle



\begin{resultbox}[Definition]
In the context of a research proposal, a problem statement is a concise, clear, and evidence-based description of a specific \textbf{researchable} issue, a \textbf{demonstrable} gap in knowledge, or a \textbf{measurable} real-world challenge that your research is \textbf{designed to solve}.
\end{resultbox}

\vspace{1em}

\textbf{\large Structured Problem Statement}\par\vspace{0.5em}

\textbf{1. What is Known}\\
The literature on Femomenología del color
 has established foundational knowledge across Across the 5 reviewed studies on Femomenología del color. Key studies have contributed empirical evidence and theoretical frameworks.

\vspace{0.5em}

\textbf{2. What is Missing}\\
Geographic scope: Most studies on Femomenología del color
 may be concentrated in specific regions, limiting generalizability to diverse populations.

\vspace{0.5em}

\textbf{3. Affected Stakeholders}\\
Researchers, practitioners, and policymakers working with Femomenología del color
 are directly affected by this gap, as are the populations ultimately served by evidence-based interventions.

\vspace{0.5em}

\textbf{4. Consequences of Inaction}\\
Without addressing this gap, decision-makers will continue to rely on incomplete or non-generalizable evidence, potentially leading to ineffective strategies and missed opportunities for improvement in Femomenología del color
.

\vspace{0.5em}

\textbf{5. How the Study Addresses the Gap}\\
The proposed study will employ rigorous methodology to directly investigate geographic scope: most studies on femomenología del color
 may be concentrated in specific regions, limiting generalizability to diverse populations., generating actionable evidence to fill this critical gap in the literature.

\vspace{1em}

\begin{resultbox}[Full Problem Statement]
Despite advances in understanding Femomenología del color
, a critical gap remains: geographic scope: most studies on femomenología del color
 may be concentrated in specific regions, limiting generalizability to diverse populations.. This is problematic because current evidence is insufficient to guide effective practice across diverse contexts. The proposed study addresses this gap by conducting rigorous, contextually grounded research that will produce generalizable findings and actionable recommendations for stakeholders in the field.
\end{resultbox}

\vspace{1em}

\textbf{\large Validation Checklist}\par\vspace{0.3em}

\begin{itemize}[leftmargin=*, label=$\square$]
    \item Is the problem \textbf{researchable}? (Can it be investigated through empirical methods?)
    \item Is the gap \textbf{demonstrable}? (Can you cite evidence that it exists?)
    \item Is the challenge \textbf{measurable}? (Can outcomes be quantified or assessed?)
    \item Are \textbf{stakeholders} clearly identified?
    \item Are \textbf{consequences of inaction} articulated?
    \item Does the statement tell reviewers \textbf{why}, \textbf{what}, and \textbf{how}?
\end{itemize}

\vfill
\begin{center}
\textit{Problem statement generated by Research Accelerant Agent on \today.}
\end{center}

\end{document}